\section{CONVERGENCE ANALYSIS}\label{sec:convergence-analysis}
\added[id=cuiys]{
In this section, we analyze the convergence behavior of the proposed federated learning algorithm under partial participation, dynamic weighting, and proxy-client tracking.
}
\subsection{Update Rule}
\added[id=cuiys]{
The goal is to minimize the global objective
$F(\omega) = \sum_{k=1}^K p_k F_k(\omega)$, where $p_k \ge 0$ and $\sum_{k=1}^K p_k = 1$.
At communication round $t$, a subset of clients $S_t$ is sampled; each selected client performs $E$ local SGD steps, and the resulting updates are aggregated.
}
\added[id=cuiys]{
To characterize the effects of partial participation and dynamic weighting, we decompose the practical global update into an ideal full-participation component and a deviation term.
Define
$g_t \triangleq \sum_{k=1}^K p_k \bar g_t^k$ as the ideal full-participation gradient, where $\bar g_t^k$ is the client-level average local gradient over $E$ steps.
Define
$d_t \triangleq \sum_{k \in S_t} \alpha_{k,t} \bar g_t^k$ as the aggregated gradient under dynamic weights $\alpha_{k,t}$.
The deviation is then
$b_t \triangleq d_t - g_t$.
}

\added[id=cuiys]{
The global model update can be written as
\begin{equation}
\omega_{t+1} = \omega_t - \eta_t E(g_t + b_t) + r_t \triangleq \omega_t + \Delta_t,
\end{equation}
where $\eta_t$ is the global learning rate and $r_t$ denotes the proxy-client residual term.
}
\subsection{Assumptions}
For the convergence analysis, we adopt the following standard assumptions.

\begin{assumption}[$L$-Smoothness]\label{ass:smoothness}
The global objective $F$ and all local objectives $F_k$ are differentiable and $L$-smooth.
\end{assumption}

\begin{assumption}[Bounded Local Variance]\label{ass:local-variance}
The stochastic gradients computed locally are unbiased and have bounded variance:
$\mathbb{E}\|g_{t,e}^k - \nabla F_k(\omega_{t,e}^k)\|^2 \le \sigma_l^2.$
\end{assumption}

\begin{assumption}[Bounded Data Heterogeneity]\label{ass:data-heterogeneity}
The divergence between local and global gradients is bounded such that
$\sum_{k=1}^K p_k \|\nabla F_k(\omega)\|^2 \le \kappa^2 + M\|\nabla F(\omega)\|^2,$
for some constants $\kappa^2 \ge 0$ and $M \ge 1$.
\end{assumption}

\begin{assumption}[Bias Decoupling]\label{ass:bias-decoupling}
The total deviation $b_t$ can be decomposed into a systematic bias and sampling variance:
$b_t = \mu_t + (b_t - \mu_t),$
where the expected systematic bias is bounded by $\|\mu_t\|^2 \le B_\mu^2$, and the sampling variance is bounded by
$\mathbb{E}_{S_t}\|b_t - \mu_t\|^2 \le B_V^2.$
\end{assumption}

\begin{assumption}[Residual Decay]\label{ass:residual-decay}
The proxy residual term decays sufficiently fast across rounds:
$\|r_t\|^2 \le \frac{D}{(t+1)^{2+2\rho}},$
where $D > 0$ and $\rho > 0$.
\end{assumption}

\subsection{Main Results}
\added[id=cuiys]{
Under the above assumptions, we now present the main convergence result for non-convex objectives.
}


\begin{theorem}[Convergence Bound]
\added[id=cuiys]{
Suppose Assumptions 1--5 hold. Let the learning rate be $\eta_t = \frac{\eta_0}{\sqrt{t+1}}$, with $\eta_0$ chosen such that $\eta_t E L \le c$ for a sufficiently small constant $c$. After $T$ communication rounds, the minimum expected squared gradient norm satisfies
}
\begin{align}
\min_{0 \le t < T} \mathbb{E}\|\nabla F(\omega_t)\|^2
&\le \mathcal{O}\!\left(\!\frac{\Delta_F}{E\sqrt{T}}\!\right)\!+\!\mathcal{O}\!\left(\!\frac{C_R}{E\sqrt{T}}\!\right)
\!+\!\mathcal{O}\!\left(\!B_\mu^2\!\right) \notag \\
&\!+\!\mathcal{O}\!\left(\!\frac{(\sigma_l^2\!+\!\kappa^2\!+\!B_V^2)E\ln T}{\sqrt{T}}\!\right),
\end{align}
\added[id=cuiys]{
where $\Delta_F = F(\omega_0) - F_{\inf}$, and $C_R < \infty$ is an absolute constant that bounds the infinite series induced by proxy residuals.
}
\end{theorem}

\added[id=cuiys]{
The complete proof is provided in the Appendix. Taking $T \to \infty$ yields the following asymptotic characterization.
}

\added[id=cuiys]{
\begin{corollary}[Asymptotic Error Floor]
Under the same conditions as Theorem 1, as $T \to \infty$, the algorithm converges to a neighborhood of a stationary point characterized by
\begin{equation}
\lim_{T \to \infty} \min_{0 \le t < T} \mathbb{E}\|\nabla F(\omega_t)\|^2 \le \mathcal{O}(B_\mu^2).
\end{equation}
\end{corollary}
}


\subsection{Discussion and Theoretical Insights}
\added[id=cuiys]{
Theorem 1 and Corollary 1 clarify how each error source contributes to optimization dynamics under dynamic partial participation.
By Assumption 5, the proxy residual term is summable, and its cumulative effect is bounded by a finite constant $C_R$.
After normalization by the effective progress term $\sum_t \eta_t = \Theta(\sqrt{T})$, this contribution decays as $\mathcal{O}(1/\sqrt{T})$.
}

\added[id=cuiys]{
Stochastic effects from local gradient noise ($\sigma_l^2$), data heterogeneity ($\kappa^2$), and client-sampling variance ($B_V^2$) scale with
$\sum_t \eta_t^2 = \Theta(\ln T)$.
Under the decaying stepsize schedule, these terms contribute
$\mathcal{O}\!\left(\frac{\ln T}{\sqrt{T}}\right)$
to the stationarity bound.
}

\added[id=cuiys]{
The systematic bias term $B_\mu^2$ remains non-vanishing and determines the asymptotic error floor.
Therefore, without an explicit debiasing mechanism, the method converges to a neighborhood of a stationary point rather than to an exact stationary point.
}